\section{Prometheus Architecture}
\label{sec:architecture-main}

Prometheus is a candidate architecture in which a neural controller constructs
executable operations, calls them on a persistent internal state, and retains
their acquired meanings during subsequent learning. Its computational unit is a
Block with a marked state interface, an operation bank, and learned input,
calling, and output maps. The architecture separates the representation of an
operation from the neural decision to invoke it, while training both through
terminal task feedback. Repeated Blocks are intended to communicate through
explicit state ports rather than independently classify intermediate outputs.

\subsection{State and executable operations}

A Block maintains neural parameters $\theta$, a logical state
$q\in\{1,\ldots,D\}$, and a bank of based loops
$\{\gamma_{k,\theta}\}_{k=1}^{K}$. A supplied covering
$p:\mathcal E\rightarrow X$ has a fixed control basepoint $x_*$ and a marked
$D$-element fiber. Lifting a loop transports the current marked state to another
point of that fiber, defining an endpoint permutation
$\Pi_k=\rho([\gamma_{k,\theta}])$. Consecutive calls carry the same state through
their operations. Their order therefore determines a genuine composed
computation, including when the endpoint actions do not commute.

The bank stores operation inputs, source versions, and execution receipts.
Endpoint maps are obtained from actual continuation and may subsequently be
cached as derived computational state. Such reuse distinguishes a paid cold
identification from a warm discrete application. A receipt certifying unchanged
semantics does not assert that a newly deformed trajectory has been integrated.
The numerical precision, domain checks, cache payload, and source observations
are part of the execution interface.

\subsection{Shared neural control and terminal learning}

Shared features $h_\theta(e)$ feed both a loop head, producing raw path
coordinates $a_\theta(e)$, and call heads, producing probabilities over bank
entries and an identity operation. Learned encoder and readout logits define
row-stochastic matrices $E_\theta$ and $R_\theta$. For a token $b_t$, optional
public context $c$, and $S$ call slots, the differentiable computation is
\begin{align}
 \alpha_0&=E_\theta[x,:],\\
 A_{b_t,c}&=\prod_{s=1}^{S}
       \left[\sum_{k=0}^{K}p_{\theta,s}(k\mid b_t,c)\Pi_k\right],
       \qquad \Pi_0=I,\\
 p_\theta(y\mid x,b_{1:T},c)&=
       [\alpha_0 A_{b_1,c}\cdots A_{b_T,c}R_\theta]_y.
 \label{eq:main-block-computation}
\end{align}
The matrices $\Pi_k$ are detached observations of executed loops. There is no
additional trainable permutation-score tensor. Gradients through the smooth
encoder, calls, and readout are combined with the class-boundary contribution
derived in Section~\ref{sec:monodromy-learning-main}. Because the backbone is
shared, either contribution can change several parts of the computation.

The native interface selects the encoder, calls, and readout by their hard
choices and executes the resulting program. Soft likelihood supplies learning
credit; successful acquisition additionally requires a correct native program.
The target specifies a terminal output, without supplying an intermediate
state, desired loop, or correct call. Known substrate equations and any observed
full-fiber actions are disclosed information available to comparison methods.

\subsection{Formation, retention, and extensibility}

An unadmitted entry can move between executable classes while terminal feedback
shapes its controller and binding. Admission installs the acquired operation
and its calling contract after deterministic execution and confidence checks.
Later updates preserve admitted loop actions, the complete old
input--call--readout relation, and agreement between the native and soft
interfaces over a specified length range. The reference records the acquired
mechanism rather than a teacher transition table. Subsequent compression or
reuse must preserve that contract.

The proposed general Block includes typed ports for state expansion and
auxiliary storage. Expansion requires an explicit correspondence between old
and new states and compatible old operations. An irreversible logical update
can use a reversible lift on the core and a fresh auxiliary register, with a
projection defining its visible result. The consumed blank, retained garbage,
and later cleanup are charged resources, following the established reversible
computation construction~\cite{bennett1973logical}. Local factorization permits
register storage without enumerating the entire history state space; efficient
learning and execution still require their own construction. Continuous
inputs and stacked Blocks additionally require compatible learned encodings
and a specified protected input domain.

The tested integrated component instantiates the encoder, shared controller,
calls, and loop bank in a fixed three-state covering. It provides an executable
example of the learning interface. The broader architecture specifies the
interfaces for growth, noninvertible computation, and composition across
Blocks; validating their common implementation precedes the primary benchmark.
That benchmark will freeze the complete architecture and retrain every method
under the same five-layer evaluation protocol. Resource comparisons will count
neural parameters together with auxiliary state, program and cache storage,
formation queries, rejected proposals, protection, and cold and warm execution.
